  

%%%%%%%%%%%%%%Empieza la presentacion%%%%%%%%%%%%%%%%%%%%%%%%%%%%%%%%%%%%

%\begin{frame}
%\frametitle{Esquema}
 %   \tableofcontents
%\end{frame}

\section{Sistemas LTI y Filtros Discretos}


%%%%%%%%%%%%%%%%%%%%%%%%%%%%%%

\begin{frame}
\frametitle{Sistemas LTI y Filtros Discretos I}

\footnotesize

\begin{itemize}
\item La respuesta y(n) de un sistema LTI discreto a una entrada  x(n) se calcula:
\begin{itemize}
\footnotesize
\item Como la convolución entre  $x(n)$ y la respuesta del sistema a la entrada impulso $h(n)$:
	$\alert{y(n)}\!=\!T[x(n)]\!=\!(h*x)(n)\!\alert{=\!\sum_{m=-\infty}^\infty x(m)h(n\!-\!m)}$
\vspace{0.1cm}
\item A través del producto de las transformadas de Fourier de la entrada y de la función ponderatriz: $Y(e^{jw})=H(e^{jw})X(e^{jw})$

\vspace{0.1cm}
\item DTFT: $X(e^{jw})=\sum_{n=-\infty}^{\infty} x(n) e^{-jwn}$
\end{itemize}


\end{itemize}

\alert{¿Podemos usar el producto de DFTs para calcular la respuesta de sistemas LTI ?}


\begin{itemize}

\visible<2->{

\item Las operaciones correspondientes en el caso de la DFT son: 

\vspace{-0.1cm}

\begin{itemize}
\footnotesize
\item $y(n)=\sum_{m=0}^{N-1} x(m)h((m-n)_N)$ para señales de longitud $\le N$
\vspace{0.1cm}

\item $Y(k)=H(k)X(k)$
\vspace{0.1cm}
\item $X(k)=\sum_{n=0}^{N-1} x(n)e^{j\frac{2\pi}{N}kn}$ para señales de longitud $\le N$
\end{itemize}

}

\visible<3->{

\item La DFT y la convolución se definen para señales de longitud finita. Por lo tanto, para poder realizar las operaciones con la DFT necesitamos que:

\begin{itemize}
\footnotesize
\item $x(n)$ sea una señal de longitud finita $L\le N$
\item $h(n)$ sea una señal de longitud finita $M \le N$  $\rightarrow $ sistema FIR.
\item Para que no aparezcan elementos del desplazamiento circular, $N \ge L+M-1$


\end{itemize}

}

\end{itemize}

\end{frame}


\begin{frame}
\frametitle{Sistemas LTI y Filtros Discretos II}

\footnotesize
\begin{itemize}
\item Para poder calcular la respuesta de un sistema a través del producto de las DFTs, necesitamos:
\begin{itemize}
\footnotesize
\item Sistema sea FIR
\item El número de muestras de las DFTs es $N>=L+M-1$, con $L$ longitud de la señal de entrada $x(n)$ y $M$ longitud de la respuesta del sistema a la función impulso $h(n)$.
\item Se rellenará con ceros las secuencias $x(n)$ y $h(n)$ para lograr que las secuencias sean de duración sea $N$ antes de calcular sus DFTs.
\end{itemize}

\visible<2->{

\item Si el sistema $h(n)$ es IIR:
\begin{itemize}
\footnotesize
\item Podemos obtener la DFT como $H(k)=H(e^{\frac{2\pi k}{N}})$, pero no utilizarlo para calcular la respuesta del sistema a cualquier entrada.
\end{itemize}}

\visible<3->{

\item Si la entrada $x(n)$ no tiene longitud finita:
\begin{itemize}
\footnotesize
\item Podemos obtener la DFT como $X(k)=X(e^{\frac{2\pi k}{N}})$, pero no utilizar la transformada para calcular la respuesta de cualquier sistema LTI a dicha entrada.
\end{itemize}
 }
 
 \visible<4->{
 
\item Las entradas $x(n)$ y funciones ponderatrices $h(n)$ periódicas son un caso particular de señales de longitud finita, considerando como longitud finita de ambas señales su periodo correspondiente.
}
 
\end{itemize}

\end{frame}


%%%%%%%%%%%%%%%%%%%%%%%%%%%%%%%

\begin{frame}
\frametitle{Ejemplo de  respuesta sistema LTI}

\footnotesize
\begin{itemize}
\item  \color{olivegreen} Calcular la respuesta del sistema $h(n)=\{0,2\}$ a la entrada  $x(n)=\{1,3,4\}$.
\color{gray}
\item A traves de la convolución:

$\left . \!\!\!\!\! \begin{array}{l} y(0)=\sum_{n=0}^{1}h(n)x(0-n)=0\cdot 1 + 2 \cdot 0=0\\
y(1)=\sum_{n=0}^{1}h(n)x(1-n)=0\cdot 3 + 2 \cdot 1=2\\
y(2)=\sum_{n=0}^{1}h(n)x(2-n)=0\cdot 4+ 2 \cdot 3=6\\
y(3)=\sum_{n=0}^{1}h(n)x(3-n)=0\cdot 0 + 2 \cdot 4=8\end{array}\!\!\right \}$ $\rightarrow y(n)=\{0,2,6,8\}$
\item A través de la DFT:

$N=3+2-1=4$

$x(n)=\{1,3,4,0\}$ 

$\left . \!\!\!\!\! \begin{array}{l} 
X(0)\!=\!\sum_{n=0}^{N-1}\!x(n)e^{-j\frac{2\pi}{N}0n}\!=\!1\!+\!3\!+\!4\!=\!8\\
X(1)\!=\!\sum_{n=0}^{N-1}\!x(n)e^{-j\frac{2\pi}{N}1n}\!=\!1\!-\!3j\!-\!4\!=\!-\!3\!-\!3j\\
X(2)\!=\!\sum_{n=0}^{N-1}\!x(n)e^{-j\frac{2\pi}{N}2n}\!=\!1\!-\!3\!+\!4\!=\!2\\
X(3)\!=\!\sum_{n=0}^{N-1}\!x(n)e^{-j\frac{2\pi}{N}3n}\!=\!1\!+\!3j\!-\!4\!=\!-3\!+\!3j\end{array}\!\!\right \}$ $\rightarrow X(k)\!=\!\{8,\!-3\!-\!3j,2,-3\!-3j\}$



$h(n)=\{0,2,0,0\}\:\:\:\:\:\:\:\:\:\:\:\:\:\:\:\:\:\:\:\:\:\:\:\:\:\:\:\:\:\:\:\:\:\:\:\:\:\:\:\:\:\:\:\:\:\:\:\:\:$ $\rightarrow H(k)=\{2,-2j,2,2j\}$

$Y(k)\!=\!X(k)H(k)\!=\!\{16,-6\!+\!6j,4,-6\!-\!6j\}$ $\:\rightarrow y(n)=\{0,2,6,8\}$


\end{itemize}

\end{frame}
